%% kapa: exam=2 round=26 year=2015 period=1 subject=appraisal_practice subject_ko=감정평가실무 seq=2 number=2 points=30
%% source: Q-Net 감정평가사 기출문제 (cst00309) artlSeq=5149601 file=extracted/1008/5149601/2015년 제26회 감정평가사 2차시험 기출문제/1교시(감정평가실무).pdf
%% conversion: kordoc 4.24.51 -> tools/gen2.py (auto); review=proofread-gpt-sol; flags=none
\begin{longquestion}{【문제 2】}
감정평가사 김공정은 택지개발사업과 관련하여 보상목적의 감정평가를 의뢰받았다. 관련법규 및 이론을 참작하고 제시된 자료를 활용하여 다음의 물음에 답하시오. (30점) \par
\medskip\noindent 물음 1) 자료와 같은 내용의 구분지상권이 설정된 토지가 공익사업에 편입되어 해당 송전선을 철거하는 경우 보상목적의 구분지상권 감정평가방법에 대해 구체적으로 기술하되 각 방법의 장점과 단점도 포함하여 기술하시오. (10점) \par
\medskip\noindent 물음 2) 주어진 자료를 활용하여 대상물건의 보상액을 구하되, 적용가능한 방법을 모두 활용한 후 시산가액의 조정을 통해 구하시오. (20점) \par
<자료 1> 공익사업에 관한 사항 \par
1. 사업명: ○○지구 택지개발사업(공익사업 근거법: 택지개발촉진법) \par
2. 사업지구면적: 180,000 m\textsuperscript{2} \par
3. 사업시행자: G지방공사 \par
4. 사업추진일정 \par
1) 택지개발지구지정ㆍ고시일: 2011. 09. 09 \par
2) 보상계획공고일: 2012. 02. 20 \par
3) 실시계획승인ㆍ고시일: 2013. 04. 04 \par
<자료 2> 감정평가의 기본적 사항 \par
1. 대상물건: 경기도 A시 B읍 C리 1번지의 구분지상권 \par
2. 감정평가목적: 보상 \par
3. 기준시점: 2015. 09. 02 \par
<자료 3> 구분지상권에 관한 사항 \par
1. 구분지상권자: H전력공사 \par
2. 구분지상권의 목적: 154 kV 가공 송전선로 건설 \par
※ 가공 송전선로: 송전철탑을 통해 공중으로 설치한 송전선로 \par
3. 구분지상권의 범위: 경기도 A시 B읍 C리 1번지 토지 상공 15 m에서 30 m까지의 공중공간(선하지면적: 300 m\textsuperscript{2}) \par
※ 선하지면적은 구분지상권 설정면적을 말함 \par
4. 구분지상권의 존속기간: 해당 송전선로 존속시까지 \par
5. 구분지상권 설정일: 2010. 03. 03 \par
6. 당시 보상액(구분지상권 설정대가): 32,000,000원 \par
7. 특약사항: 존속기간 동안 구분지상권 설정대가의 증감은 없음 \par
8. 기타사항: 송전선로가 필지의 중앙을 통과함 \par
<자료 4> 토지에 관한 사항 \par
1. 소재지: 경기도 A시 B읍 C리 1번지 \par
2. 면적: 300 m\textsuperscript{2}, 지목: 전, 실제 이용상황: 전 \par
3. 접면도로: 폭 6 m의 도로와 접함 \par
4. 토지이용계획의 변동사항 \par
1) 2010. 01. 01 ～ 2013. 04. 03: 자연녹지지역 \par
2) 2013. 04. 04 ～ 2015. 09. 02: 일반주거지역(택지개발사업으로 인해 변경) \par
<자료 5> 주변지역 현황 \par
1. 구분지상권이 설정된 토지(경기도 A시 B읍 C리 1번지 토지) 주변은 송전선 건설 당시 농지지대에서 주택지대로 전환되는 중이었고(단독주택이 지속적으로 건설되고 있었음), 당시 인근지역에 속한 토지로서 ○○택지개발지구 인근의 지구 밖 토지 대부분은 기준시점 현재 단독주택부지로 이용하고 있음 \par
2. 조사결과 기준시점 현재 ○○택지개발지구와 접한 지구 밖 토지의 표준적이용은 2층의 단독주택부지이며, 주택의 표준적인 각 층의 층고는 3.5 m임 \par
<자료 6> 표준지 공시지가 자료 \par
{\small\centering
\resizebox{\ifdim\width>\linewidth\linewidth\else\width\fi}{!}{%
\begin{tabular}{|c|c|c|c|c|c|c|}
\hline
\makecell{기호} & \makecell{소재지 지번} & \makecell{면적\\(m\textsuperscript{2})} & \makecell{지목} & \makecell{이용상황} & \makecell{도로교통} & \makecell{형상지세} \\
\hline
\makecell{가} & \makecell{B읍 C리 10} & \makecell{360} & \makecell{답} & \makecell{과수원} & \makecell{세로(불)} & \makecell{세장형 완경사} \\
\hline
\makecell{나} & \makecell{B읍 C리 250} & \makecell{280} & \makecell{대} & \makecell{주거나지} & \makecell{세로(가)} & \makecell{세장형 평지} \\
\hline
\makecell{다} & \makecell{B읍 D리 500} & \makecell{350} & \makecell{전} & \makecell{과수원} & \makecell{세로(가)} & \makecell{가장형 평지} \\
\hline
\end{tabular}}\par}

※ 표준지 ‘가’ ～ ‘다’는 모두 동일수급권내에 소재함 \par
(계속) \par
{\small\centering
\resizebox{\ifdim\width>\linewidth\linewidth\else\width\fi}{!}{%
\begin{tabular}{|c|c|c|c|c|}
\hline
\multirow{2}{*}{\makecell{기호}} & \multirow{2}{*}{\makecell{용도지역}} & \multicolumn{2}{|c|}{\makecell{공시지가(원/m\textsuperscript{2})}} & \multirow{2}{*}{\makecell{비고}} \\
\cline{3-4}
 &  & \makecell{공시기준일} & \makecell{공시지가} &  \\
\hline
\multirow{3}{*}{\makecell{가}} & \multirow{2}{*}{\makecell{자연녹지}} & \makecell{2011.01.01} & \makecell{300,000} & \multirow{3}{*}{\makecell{○○지구 택지개발사업지구 내의 토지로\\도시계획시설 도로에 40 \% 저촉함}} \\
\cline{3-4}
 &  & \makecell{2013.01.01} & \makecell{380,000} &  \\
\cline{2-4}
 & \makecell{일반주거} & \makecell{2015.01.01} & \makecell{480,000} &  \\
\hline
\multirow{3}{*}{\makecell{나}} & \multirow{2}{*}{\makecell{자연녹지}} & \makecell{2011.01.01} & \makecell{340,000} & \multirow{3}{*}{\makecell{○○지구 택지개발사업지구 내의 토지로 가공 송전선으로 인해 구분지상권이\\설정되어 있음}} \\
\cline{3-4}
 &  & \makecell{2013.01.01} & \makecell{420,000} &  \\
\cline{2-4}
 & \makecell{일반주거} & \makecell{2015.01.01} & \makecell{500,000} &  \\
\hline
\multirow{3}{*}{\makecell{다}} & \multirow{3}{*}{\makecell{자연녹지}} & \makecell{2011.01.01} & \makecell{300,000} & \multirow{3}{*}{\makecell{○○지구 택지개발사업지구 밖의 토지로 가공 송전선으로 인해 구분지상권이\\설정되어 있음}} \\
\cline{3-4}
 &  & \makecell{2013.01.01} & \makecell{360,000} &  \\
\cline{3-4}
 &  & \makecell{2015.01.01} & \makecell{500,000} &  \\
\hline
\end{tabular}}\par}

※ 도시계획시설 도로에 저촉하는 표준지의 경우 해당부분에 대해 20 \%의 감가율을 적용하여 공시하였음 \par
<자료 7> 시점수정 자료: 경기도 A시 지가변동률 \par
{\small\centering
\resizebox{\ifdim\width>\linewidth\linewidth\else\width\fi}{!}{%
\begin{tabular}{|c|c|c|}
\hline
\multirow{2}{*}{\makecell{기간}} & \multicolumn{2}{|c|}{\makecell{지가변동률(단위: \%)}} \\
\cline{2-3}
 & \makecell{주거지역} & \makecell{녹지지역} \\
\hline
\makecell{2010. 01. 01 ～ 2010. 12. 31(누계)} & \makecell{1.103} & \makecell{2.758} \\
\hline
\makecell{2011. 01. 01 ～ 2011. 12. 31(누계)} & \makecell{2.154} & \makecell{3.085} \\
\hline
\makecell{2012. 01. 01 ～ 2012. 12. 31(누계)} & \makecell{2.060} & \makecell{2.072} \\
\hline
\makecell{2013. 01. 01 ～ 2013. 12. 31(누계)} & \makecell{2.058} & \makecell{2.085} \\
\hline
\makecell{2014. 01. 01 ～ 2014. 12. 31(누계)} & \makecell{2.064} & \makecell{3.082} \\
\hline
\makecell{2015. 01. 01 ～ 2015. 07. 31(누계)} & \makecell{-0.130} & \makecell{-0.120} \\
\hline
\makecell{2015. 07. 01 ～ 2015. 07. 31(당월)} & \makecell{0.060} & \makecell{0.072} \\
\hline
\end{tabular}}\par}

※ 2015년 8월 이후 지가변동률은 미고시 상태이며 2015년 7월 지가변동률과 동일하게 변동하는 것으로 추정함 \par
<자료 8> 보상사례 자료 \par
1. 보상물건: 경기도 A시 B읍 D리 500번지[<자료 6>의 기호 ‘다’ 토지임]의 구분지상권(구분지상권 설정일: 2010. 03. 03) \par
2. 보상사유: 도시개발사업(사업인정고시일: 2013. 05. 01, 사업지구면적: 100,000 m\textsuperscript{2})에 편입되어 154 kV 가공 송전선로 철거 \par
3. 보상액: 37,400,000원 \par
4. 보상액 감정평가시 기준시점: 2015. 07. 31 \par
5. 구분지상권의 범위: 경기도 A시 B읍 D리 500번지 토지 상공 16 m에서 30 m까지의 공중공간(선하지면적: 280 m\textsuperscript{2}) \par
6. 구분지상권의 존속기간: 해당 송전선로 존속시까지 \par
7. 특약사항: 존속기간 동안 구분지상권 설정대가의 증감은 없음 \par
8. 기타사항: 송전선로가 필지의 중앙을 통과함 \par
<자료 9> 구분지상권의 가치형성요인 비교 자료 등 \par
1. 조사결과 가공 송전선로를 위한 구분지상권의 가치는 해당 부지의 지역요인 및 개별요인, 송전선로로 인한 입체이용저해율 및 추가보정률(쾌적성 저해요인, 시장성 저해요인, 기타 저해요인), 선하지면적에 영향을 받음 \par
2. 조사결과 <자료 8>의 보상사례와 대상물건은 <자료 10> 및 <자료 11>의 내용과 같이 비교치가 산정됨 \par
3. 조사결과 구분지상권의 가치는 지가변동률과 동일하게 변동함 \par
<자료 10> 지가형성요인 비교 자료 \par
1. 지역요인 비교치: 표준지와 비교한 B읍 C리 1번지의 비교치 \par
{\small\centering
\resizebox{\ifdim\width>\linewidth\linewidth\else\width\fi}{!}{%
\begin{tabular}{|c|c|c|}
\hline
\makecell{표준지 ‘가’} & \makecell{표준지 ‘나’} & \makecell{표준지 ‘다’} \\
\hline
\makecell{1.00} & \makecell{1.00} & \makecell{1.10} \\
\hline
\end{tabular}}\par}

※ 비교치는 표준지공시지가의 공시기준일이 상이해도 동일하게 적용함 \par
2. 개별요인 비교치: 표준지와 비교한 B읍 C리 1번지의 비교치 \par
{\small\centering
\resizebox{\ifdim\width>\linewidth\linewidth\else\width\fi}{!}{%
\begin{tabular}{|c|c|c|}
\hline
\makecell{표준지 ‘가’} & \makecell{표준지 ‘나’} & \makecell{표준지 ‘다’} \\
\hline
\makecell{1.10} & \makecell{0.90} & \makecell{1.00} \\
\hline
\end{tabular}}\par}

※ 비교치는 표준지공시지가의 공시기준일이 상이해도 동일하게 적용함 \par
3. 그 밖의 요인 보정치 \par
토지의 감정평가에 적용할 그 밖의 요인 보정치는 공시기준일에 상관없이 표준지 ‘가’ ～ ‘다’ 모두 1.30으로 적용함 \par
<자료 11> 보정률 산정 자료 \par
1. 건조물의 이격거리 \par
건조물은 가공전선의 전압 35 kV 이하는 3 m, 35 kV를 초과하는 경우에는 초과하는 10 kV 또는 그 단수마다 15 cm를 가산한 수치씩 이격하여야 함 \par
2. 주택지대의 층별효용지수 1층: 100, 2층: 100 \par
3. 입체이용률 배분표 \par
{\small\centering
\resizebox{\ifdim\width>\linewidth\linewidth\else\width\fi}{!}{%
\begin{tabular}{|c|c|c|c|c|}
\hline
\makecell{구분} & \makecell{건물이용률(α)} & \makecell{지하이용률(β)} & \makecell{그 밖의 이용률(γ)} & \makecell{γ의 상하 배분비율} \\
\hline
\makecell{주택지대ㆍ\\택지후보지지대} & \makecell{0.7} & \makecell{0.15} & \makecell{0.15} & \makecell{3:1} \\
\hline
\makecell{농지지대} & \makecell{0.8} & \makecell{0.1} & \makecell{0.1} & \makecell{4:1} \\
\hline
\end{tabular}}\par}

4. 추가보정률 산정기준표 \par
{\small\centering
\resizebox{\ifdim\width>\linewidth\linewidth\else\width\fi}{!}{%
\begin{tabular}{|c|c|c|c|}
\hline
\multirow{2}{*}{\makecell{구분}} & \multicolumn{2}{|c|}{\makecell{추가보정률 적용범위}} & \multirow{2}{*}{\makecell{상ㆍ중ㆍ하 구분 기준}} \\
\cline{2-3}
 & \makecell{주택지대ㆍ택지후보지지대} & \makecell{농지지대} &  \\
\hline
\makecell{쾌적성\\저해요인} & \makecell{상: 10.0 \%\\중: 7.5 \%\\하: 5.0 \%} & \makecell{상: 5.0 \%\\중: 4.0 \%\\하: 3.0 \%} & \makecell{송전선로의 높이를 기준으로 구분 적용\\ㆍ10 m이하: 전압에 관계없이 ‘상’\\ㆍ10 m 초과 20 m이하: 154 kV이하는 ‘중’,\\154 kV 초과는 ‘상’\\ㆍ20 m 초과: 754 kV이상은 ‘상’, 345 kV\\이상은 ‘중’, 154 kV 이하는 ‘하’} \\
\hline
\makecell{시장성\\저해요인} & \makecell{상: 10.0 \%\\중: 7.0 \%\\하: 4.0 \%} & \makecell{상: 7.0 \%\\중: 5.0 \%\\하: 3.0 \%} & \makecell{선하지면적비율 또는 송전선로의 통과\\위치를 기준으로 구분 적용\\ㆍ선하지면적비율이 40 \%를 초과하거나\\송전선로가 필지의 중앙을 통과하는\\경우: ‘상’\\ㆍ선하지면적비율이 20 \%를 초과하거나\\송전선로가 필지의 측면을 통과하는\\경우: ‘중’\\ㆍ선하지면적비율이 20 \% 이하이거나\\송전선로가 필지의 모서리를 통과하는\\경우: ‘하’} \\
\hline
\makecell{기타\\저해요인} & \makecell{상: 10.0 \%\\중: 6.0 \%\\하: 3.0 \%} & \makecell{상: 8.0 \%\\중: 5.0 \%\\하: 3.0 \%} & \makecell{송전선로의 존속기간을 기준으로 구분 적용\\ㆍ존속기간이 30년을 초과하는 경우: ‘상’\\ㆍ존속기간이 10년을 초과하는 경우: ‘중’\\ㆍ존속기간이 10년 이하인 경우: ‘하’} \\
\hline
\multicolumn{4}{|c|}{\makecell{추가보정률 산정 기준: 각 해당 항목을 가산하여 산정}} \\
\hline
\end{tabular}}\par}

% [원문 그대로] 추가보정률 표의 ‘754 kV’는 원문 수치를 보존함.
<자료 12> 기타 자료 \par
1. 일시금운용이율(또는 환원율): 연 5.0 \% \par
\end{longquestion}
